% Drawing 138/06-1 -- folder 138, batch 06, archive page 103 (PDF page 104).
% Transcription: un segment, $\bullet$ en $0$ (mult.\ $1$), $\times$ en $1$ (mult.\ $1$) ; $\infty$ de multiplicité $1$.
% Kind: dessin d'enfant (one edge). Method 3 (encoded, rendered from the data). Status: ok.
% Notes: A segment from a bullet at 0 (multiplicity 1) to a cross at 1 (multiplicity 1); infinity, multiplicity 1, apart. Encoded: the trivial dessin (degree 1). Notation of the page: a bullet over the value 0 is a black vertex, a cross over 1 a white vertex, a small circle over infinity the pole in a face; the number under each symbol is the multiplicity.
% Coordinates: 1 unit = 1 mm = 1/100 of the crop width (bbox in page fractions: [0.484, 0.073, 0.676, 0.129]).
% Redrawn by Claude (claude-opus-5-5), issue #73, 2026-11-07; not checked by a human; not compiled (no TeX in the session).
% Rights: reproduces part of the work -- behind RENDER_TRANSCRIPTS, never published by itself.
\documentclass[tikz,border=3pt]{standalone}
\usepackage{amsmath,amssymb}
\usetikzlibrary{patterns,arrows.meta}
\tikzset{pencil/.style={black!45}, struck/.style={black!35}, mark/.style={circle,draw,inner sep=1.3pt}, pt/.style={circle,fill,inner sep=0.9pt}, opt/.style={circle,draw,fill=white,inner sep=1.1pt}}
\begin{document}
\begin{tikzpicture}[x=1mm,y=1mm,line cap=round,line join=round,font=\scriptsize]
\draw[very thick] (13.5,-21.2) -- (38.5,-21.2);
\node[pt] at (13.5,-21.2) {};
\draw (36.3,-19) -- (40.6,-23.3);
\draw (36.3,-23.3) -- (40.6,-19);
\node[inner sep=1pt] at (13.5,-12.5) {\tiny $0$};
\node[inner sep=1pt] at (13.5,-29.8) {\tiny $1$};
\node[inner sep=1pt] at (38.5,-12.5) {\tiny $1$};
\node[inner sep=1pt] at (38.5,-29.8) {\tiny $1$};
\node[inner sep=1pt] at (86.5,-16.3) {\tiny $\infty$};
\node[opt] at (86.5,-25) {};
\node[inner sep=1pt] at (86.5,-33.7) {\tiny $1$};
\end{tikzpicture}
\end{document}
